\documentclass[10pt,a4paper]{amsart}
\usepackage[margin=19mm]{geometry}
\usepackage{amsmath,amssymb,mathtools}
\usepackage[scr=rsfs,cal=euler]{mathalfa}
\usepackage{xfrac}
\usepackage[unicode,hidelinks,bookmarks=false]{hyperref}
\newcommand{\ie}{i.e.\ }
\newcommand{\cf}{cf.\ }
\newcommand{\resp}{respectively\ }
\newcommand{\Subsection}{\subsection}
\providecommand{\nobreakdash}{\nobreak\hbox{-}}
\setlength{\emergencystretch}{2em}
\multlinegap0pt
\usepackage{bm}
\usepackage{bbm}
\usepackage{dsfont}
\usepackage{xspace}
\usepackage{cancel}
\usepackage{mathtools}
\usepackage{amsthm}
\makeatletter
\@ifundefined{lemma}{\newtheorem{lemma}{Lemma}}{}
\@ifundefined{theorem}{\newtheorem{theorem}{Theorem}}{}
\@ifundefined{proposition}{\newtheorem{proposition}[theorem]{Proposition}}{}
\makeatother
\title[Epi-convergence in distribution of normal integrands with ap]{Epi-convergence in distribution of normal integrands with applications to sets of epsilon-optimal solutions}
\author{Dietmar Ferger}
\date{Source: arXiv:2507.16297v1 (CC BY 4.0)}
\begin{document}
\maketitle

\noindent\emph{Source and licence.} Epi-convergence in distribution of normal integrands with applications to sets of epsilon-optimal solutions. Version arXiv:2507.16297v1 (CC BY 4.0), distributed under the Creative Commons Attribution 4.0 International licence, as recorded in the arXiv licence metadata for this exact version and shown on the abstract page https://arxiv.org/abs/2507.16297v1. Full rights notice: licenses/arxiv-2507.16297-CC-BY-4.0.txt.\par\smallskip

\noindent\textit{Editorial scope.} Counted unit: the Theorem whose source label is \texttt{sufficient} in the article (source0.tex lines 431--440, source label \texttt{sufficient}), delivered together with the article's own complete original proof of it (source0.tex lines 442--483, one \texttt{proof} environment of 42 lines). No mathematics is written, repaired or re-derived: statement and proof are the article's own text line for line. Required context retained uncounted: phi (lines 222--224); UD (lines 275--283); necessary (lines 398--408); limgreater (lines 399--402); limlessthanorequal (lines 399--402); closedballisclosedrectangle (lines 731--733); continuityset (lines 736--741); epi1 (lines 749--751). Every definition, display and label of those lines is retained unchanged. Omitted from this package and not redistributed: 1--221, 225--274, 284--397, 403--430, 441--441, 484--730, 734--735, 742--748, 752--843. Nothing inside a retained excerpt is omitted or altered: every retained body is delivered exactly as the article's lines are written, so no omission receipt of any kind accompanies this package. Citations: the retained excerpts carry no citation command at all, so no bibliography entry of the article is transcribed and no reference is redistributed. Reference adaptation: none was needed and none was made; every reference inside the delivered unit resolves inside it, so no unresolved reference or citation is produced. Printed slips kept literal: no printed word, symbol, hypothesis or number of the counted statement or of its proof was altered. Layout and class: the article's own class is \texttt{sn-jnl} with packages including graphicx, multirow, amsmath, amssymb, amsfonts, amsthm, mathrsfs, appendix, xcolor, textcomp, manyfoot, booktabs, algorithm, algorithmicx; the wrapper is the lane's pinned \texttt{amsart} editorial wrapper at 10pt on A4, which restates the theorem-like environments the retained text needs and adds the article's own macro definitions that the retained text needs (none), with extra wrapper packages (bm, bbm, dsfont, xspace, cancel, mathtools). The delivered numbering is the unit's own. Assets: no retained span contains a figure environment, a graphics-inclusion command, a PDF-page inclusion, a table or a TikZ picture; no image file is copied into this package, the package declares no asset dependency and no image file is redistributed with it. Licence: version arXiv:2507.16297v1 of this article is distributed under the Creative Commons Attribution 4.0 International licence, as recorded in the arXiv licence metadata for this exact version and shown on the abstract page https://arxiv.org/abs/2507.16297v1. A full rights notice is delivered at licenses/arxiv-2507.16297-CC-BY-4.0.txt. Characters: every retained excerpt is pure ASCII, so the delivered engine, XeLaTeX with its default Latin Modern text font, reports no missing character.
\begin{equation} \label{phi}
\phi:(S,\tau_e) \rightarrow (\mathcal{E},\sigma) \text{ defined by } \phi(f)=\text{epi}(f),
\end{equation}

\begin{theorem} \label{UD} The following two statements are equivalent:
\begin{itemize}
\item[(1)] $Q_\alpha \longrightarrow_w Q \quad \text{on } (\mathcal{F},\tau_F).$
\item[(2)]
\begin{equation} \label{convergencedeterminingclass}
 \lim_\alpha Q_\alpha(\mathcal{H}(U))=Q(\mathcal{H}(U)) \quad \text{for all }  U \in \mathcal{U}(D).
\end{equation}
\end{itemize}
\end{theorem}

\begin{proposition} \label{necessary} Weak convergence $P_\alpha \rightarrow_w P$ in $(S,\tau_e)$ entails
\begin{align}
 \lim_\alpha P_\alpha(\cap_{i=1}^m \{I_{K_j} \le a_j\}) = P(\cap_{i=1}^m \{I_{K_j} \le a_j\}), \label{limlessthanorequal} \\
 \lim_\alpha P_\alpha(\cap_{i=1}^m \{I_{K_j} > a_j\}) = P(\cap_{i=1}^m \{I_{K_j} > a_j\}). \label{limgreater}
\end{align}
Here, (\ref{limlessthanorequal}) and (\ref{limgreater}), respectively, hold for all $m \in \mathbb{N}, a_1,\ldots,a_m \in \mathbb{R}$ and
$K_1,\ldots,K_m \in \mathcal{K}$ satisfying the continuity condition
$$
  P(I_{K_j} \le a_j)=P(I_{K_j^0} < a_j), j=1,\ldots,m.
$$
\end{proposition}

\begin{equation} \label{closedballisclosedrectangle}
\overline{B}_{d \times u}((x,\alpha),r)= \overline{B}_d(x,r) \times [\alpha-r,\alpha+r].
\end{equation}

\begin{lemma} \label{continuityset} Let $P$ be a probability measure on $(S,\mathcal{B}_e)$ and $\phi$ be the homeomorphism (\ref{phi}). If $Q:= P \circ \phi^{-1}$, then for each $(x,\alpha) \in E \times \mathbb{R}$ and $r>0$ the pertaining closed ball
$\overline{B}_{d \times u}((x,\alpha),r)$ is a $Q$-continuity set, if and only if
$$
 P( I_{\overline{B}_d(x,r)} \le r+\alpha) = P( I_{(\overline{B}_d(x,r))^0} < r+\alpha).
$$
\end{lemma}

\begin{equation} \label{epi1}
 \text{epi}(f) \cap (\overline{B}_d(x,r) \times [\alpha-r,\alpha+r]) \neq \emptyset \; \Leftrightarrow \; I_{\overline{B}_d(x,r)}(f) \le r+\alpha.
\end{equation}

\begin{theorem} \label{sufficient} Let $D:=D(E_0 \times \mathbb{Q},P \circ \phi^{-1})$. Then the following statements (1)-(3) are equivalent:
\begin{itemize}
\item[(1)] $P_\alpha \rightarrow_w P$ in $(S,\tau_e)$
\item[(2)] $\lim_\alpha P_\alpha(\cap_{j=1}^m \{I_{\overline{B}(x_j,r_j)} \le r_j+\alpha_j\}) = P(\cap_{j=1}^m \{I_{(\overline{B}(x_j,r_j))^0} \le r_j+ \alpha_j\}).$
\item[(3)] $\lim_\alpha P_\alpha(\cap_{j=1}^m \{I_{\overline{B}(x_j,r_j)} > r_j+\alpha_j\}) = P(\cap_{j=1}^m \{I_{(\overline{B}(x_j,r_j))^0} > r_j+ \alpha_j\}).$
\end{itemize}
Here, the equalities in (2) and (3), respectively, hold for all
$m \in \mathbb{N}, x_1,\ldots,x_m \in E_0, r_1,\ldots,r_m \in D, \alpha_1,\ldots,\alpha_m \in \mathbb{Q}.$ (Notice that the closed balls
$\overline{B}(x_j,r_j)$ satisfy the continuity condition (\ref{continuitycondition}) below.)
\end{theorem}

\begin{proof} The necessity of (2) or (3), respectively, for weak convergence (1) follows from Proposition \ref{necessary}, because closed balls are compact and $r_j \in D$ means that $\overline{B}_{d \times u}((x_j,\alpha_j),r_j)$ is a $P \circ \phi^{-1}$-continuity set, which by Lemma \ref{continuityset} in the appendix is equivalent to
\begin{equation} \label{continuitycondition}
P(I_{\overline{B}(x_j,r_j)} \le r_j+\alpha_j )= P(I_{(\overline{B}(x_j,r_j))^0} < r_j+\alpha_j).
\end{equation}
As to sufficiency we will see that it is enough to show that
\begin{equation} \label{Pphiweak}
P_\alpha \circ \phi^{-1} \rightarrow_w P \circ \phi^{-1} \quad \text{on } (\mathcal{F}(E \times \mathbb{R}),\tau_F(E \times \mathbb{R})).
\end{equation}
So, we are dealing with $(E \times \mathbb{R}, d \times u)$ as carrier space instead of $(E,d)$. (The product-metric is specified in the Appendix below.) Assume that (2) holds. Observe that
$$
 \mathcal{U}(D) = \{\bigcup_{i=1}^m \overline{B}_{d \times u}((x_i,\alpha_i),r_i): m \in \mathbb{N}, (x_i,\alpha_i) \in E_0 \times \mathbb{Q}, r_i \in D, i=1,\ldots,m\}.
$$
Let $U=\bigcup_{i=1}^m B_i \in \mathcal{U}(D)$, i.e. each $B_i$ is equal to the closed ball $\overline{B}_{d \times u}((x_i,\alpha_i),r_i), i=1,\ldots,m.$ Put $Q_\alpha := P_\alpha \circ \phi^{-1}$ and $Q := P \circ \phi^{-1}$. The \emph{inclusion-exclusion formula} yields:
\begin{equation} \label{ief}
Q_\alpha(\mathcal{H}(U))=Q_\alpha(\cup_{i=1}^m(\mathcal{H}(B_i)))=\sum_{k=1}^m (-1)^{k+1} \sum_{1 \le i_1<\ldots<i_k\le m} Q_\alpha(\cap_{s=1}^k \mathcal{H}(B_{i_s})).
\end{equation}
For each summand on the right side of equation (\ref{ief}) it follows that:
\begin{eqnarray*}
Q_\alpha(\cap_{s=1}^k \mathcal{H}(B_{i_s}))&=&P_\alpha(\cap_{s=1}^k \{f \in S: \text{epi}(f) \in \mathcal{H}(B_{i_s})\})\\
&=&P_\alpha(\cap_{s=1}^k \{f \in S: \text{epi}(f) \cap B_{i_s} \neq \emptyset\})\\
&=&P_\alpha(\cap_{s=1}^k \{I_{\overline{B}_d(x_{i_s},r_{i_s})} \le r_{i_s}+\alpha_{i_s}\}) \quad \text{by } (\ref{closedballisclosedrectangle}) \text{ and } (\ref{epi1}).
\end{eqnarray*}
Consequently by (2), $$Q_\alpha(\cap_{s=1}^k \mathcal{H}(B_{i_s})) \rightarrow P(\cap_{s=1}^k \{I_{\overline{B}_d(x_{i_s},r_{i_s})} \le r_{i_s}+\alpha_{i_s}\})=Q(\cap_{s=1}^k \mathcal{H}(B_{i_s}))$$
upon noticing (\ref{closedballisclosedrectangle}) and (\ref{epi1}) again.
With (\ref{ief}) we obtain that $Q_\alpha(\mathcal{H}(U)) \rightarrow Q(\mathcal{H}(U))$ for all $U \in \mathcal{U}(D).$ Thus the weak convergence (\ref{Pphiweak}) follows from Theorem \ref{UD}. Deduce from (\ref{Pphiweak}) that
$$
 P_\alpha \circ \phi^{-1} \rightarrow_w P \circ \phi^{-1} \quad \text{on the subspace } (\mathcal{E},\sigma).
$$
This yields $P_\alpha \rightarrow_w P$ by the Continuous Mapping Theorem, because
$$P_\alpha=(P_\alpha \circ \phi^{-1})\circ (\phi^{-1})^{-1}$$
and $\phi^{-1}:(\mathcal{E},\sigma) \rightarrow (S,\tau_e)$
is continuous.

Finally, assume that (3) holds. Then
\begin{eqnarray*}
Q_\alpha(\mathcal{M}(U))&=& Q_\alpha(\mathcal{M}(\cup_{i=1}^m B_i))= Q_\alpha(\cap_{i=1}^m \mathcal{M}(B_i))\\
                        &=&P_\alpha(\cap_{i=1}^m \{f \in S: \text{epi}(f) \in \mathcal{M}(B_i)\})\\
                        &=&P_\alpha(\cap_{i=1}^m \{f \in S: \text{epi}(f) \cap B_i = \emptyset\})\\
                        &=&P_\alpha(\cap_{i=1}^m \{I_{\overline{B}_d(x_i,r_i)} > r_i+\alpha_i\}) \quad \text{by } (\ref{closedballisclosedrectangle}) \text{ and } (\ref{epi1}).
\end{eqnarray*}
So, weak convergence (\ref{Pphiweak}) follows from Theorem \ref{UD} by complementation, which as shown above results in (1).
\end{proof}

\end{document}
